\documentclass[11pt]{article}

\input{0-parameters.tex}

\begin{document}

\begin{abstract}
    Under consistent input--output national accounting (with fixed shares $a_{ji}$), partial substitution of gas-based electricity by renewables in manufacturing does not alter aggregate TFP growth, regardless of whether energy flows are measured in final, useful, or primary form. 
    Using the $\zeta$-framework created in this paper, considering primary, final or useful energy does not change  TFP growth - except when primary energy is analysed within a broad energy sector that includes both electricity and fossil fuels.
Except for this one case: electrification and the switch from fossil to renewables has no impact on productivity
\[
   \hat{A}=0.
\]

In the CCC approach we are unfortunately in the ``exception'' case where a single energy sector sells both MWh of primary fossil fuel and electricity (both from renewables and fossil fuel). \textbf{This is why the best approach is to use useful energy.}
\end{abstract}




\section{Appendix: Switching from Gas to Renewable Electricity (Final, Useful, and Primary)}
\label{appendix_gas_renewables}

Energy accounting often distinguishes three stages of energy flows: \textbf{primary}, \textbf{final}, and \textbf{useful}. By convention:

\begin{itemize} \item \textbf{Primary energy} is energy found in raw fuels or natural sources (coal, crude oil, wind, etc.) \emph{before} any conversion or transformation processes. \item \textbf{Final energy} is the energy commodity \emph{delivered to end-use sectors} after refining, electricity generation, transmission, or other conversion steps. \item \textbf{Useful energy} is the \emph{portion of final energy} that effectively provides energy services (heat, mechanical work, lighting, etc.), after end-use device losses \citep{wayEmpiricallyGroundedTechnology2022}. \end{itemize}

This appendix illustrates how partial switching from gas-based electricity to renewable electricity in the manufacturing sector affects sectoral and aggregate productivity, under three energy-accounting approaches: final, useful, and primary.  We use the primal--dual total factor productivity (TFP) framework
\begin{subequations}
\begin{align}
\hat{A}_i  &\equiv  
\underbrace{\hat{X}_i  -  \sum_{j=1}^N a_{ji} \hat{X}_{ij}  -  \sum_{f=1}^F \tilde{l}_{fi} \hat{L}_{fi}}_{\text{primal}}
\label{eq:primal_appendix}\\
 &= 
\underbrace{-\left( \hat{p}_i  -  \sum_{j=1}^N a_{ji} \hat{p}_j  -  \sum_{f=1}^F \tilde{l}_{fi} \hat{w}_f \right)}_{\text{dual}}
\label{eq:dual_appendix}
\end{align}
\end{subequations}

\noindent The variables $a_{ji}$ are the intermediate-input cost shares (so $\sum_j a_{ji} + \sum_f \tilde{l}_{fi}=1$), and $\tilde{l}_{fi}$ are the labour-cost shares.

We consider that the only inputs to the energy sector, oil, or electricity, are only capital and labour. The manufacturing sector is using labour, capital and energy (either one input or electricity or electricity and oil).

\tableofcontents


\paragraph{Why the cost shares $a_{ji}$ remain constant}

As in Appendix \ref{appendix_useful_energy}, the expenditures shares of energy remains constant when considering useful, final or primary energy as the total energy sales remain constant.

\subsection{Set-up}

We can express energy three ways: primary, final or useful. $\beta$ is the conversion factor from gas to electricity, and $\alpha$ is the conversion factor between final electricity and useful energy,  both under the same constraint: $0 \leq \alpha, \beta\leq 1$. Primary is the initial state of the energy, final is the form that is actually used. So a quantity $\beta X_G$ of final gas-generated electricity corresponds to a primary quantity of energy $X_G$ of gas. There is no difference between primary and final energy for renewables. 



We consider three industries:
\begin{itemize}
\item \textbf{G}: Gas sector, extracts raw gas and can sell electricity generated from it,
\item \textbf{R}: Renewable energy, produces renewable electricity,
\item \textbf{M}: Manufacturing, needs electricity to run its processes.
\end{itemize}


One MWh of electricity produced by renewables contains 1 MWh of final and primary energy, and provides $\alpha$ MWh of useful energy. One MWh of gas, provides $\beta$ MWh of electricity, and $\alpha\beta$ MWh of useful energy.
 

We consider three scenarios where there is a small substitution of gas by renewables: a fraction $\delta$ of M's gas consumption is replaced by renewable electricity, holding M's total output constant $\hat{X}_M=0$. Let $\gamma$ be the relative increase in renewables final electricity output.
\begin{itemize}
\item \textbf{Primary}: The manufacturing sector initially buys $X_G$ MWh of gas from G (using it directly or transforming it in electricity via direct gas combustion on-site) and buys $X_R$ MWh electricity from R (We either consider two different sectors R and G in section \ref{subappendix_primary_energy} or a single energy sector E in section \ref{subappendix_primary_energy_two_sectors}).
\item \textbf{Final}: The gas sector directly sells gas-generated electricity. M buys final energy electricity $\beta X_G$ from the gas sector G and some final renewable electricity  $X_R$  from R (We either consider two different sectors R and G in section \ref{subappendix_final_energy} or a single energy sector E in section \ref{subappendix_final_energy_two_sectors})
\item \textbf{M}: The gas sector directly sells useful gas-generated electricity. Manufacturing M buys directly useful energy gas-generated electricity $\alpha \beta X_G$ from the gas sector and some useful electricity $\alpha X_R$ from the renewables sector 
(We either consider two different sectors R and G in section \ref{subappendix_useful_energy} or a single energy sector E in section \ref{subappendix_final_energy_two_sectors}).
\end{itemize}



 
 
 Table \ref{tab_summary_energy_prod} summarises the changes in supply and demand of energy.
 Important note (\ER{VERY IMPORTANT}), what we denote here as "final energy scenario" is the scenario where all the gas output is used to produce electricity, and some is switch from gas to renewables. In the CCC scenarios we have both i) gas-electricity switched to renewables electricity (covered by "Final energy scenario" in this note), ii) gas directly used by the manufacturing switch to renewable electricity (covered by "Primary energy scenario" in this note). \er{In other words, it is equivalent to consider the electrification of} 


\begin{table}[H]
\centering
\begin{tabular}{@{}c|cc|cc@{}}
\toprule
        & \multicolumn{2}{c|}{Before the switch} & \multicolumn{2}{c}{After the switch} \\ \midrule
        & Gas             & Renewable            & Gas            & Renewable           \\ \midrule
Primary &     $X_G$   &    $X_R$    &  $X_G(1-\delta)$  & $X_R+\beta\delta X_G=X_R(1+\gamma)$               \\
Final   &      $\beta X_G$              &  $X_R$                      &          $\beta X_G (1-\delta)$        &        $X_R+\beta\delta X_G=X_R(1+\gamma)$               \\
Useful&  $\alpha\beta X_G$              & $\alpha X_R$                      &          $\alpha\beta X_G (1-\delta)$        &        $\alpha(X_R+\beta\delta X_G)=\alpha X_R(1+\gamma)$                   \\ \bottomrule
\end{tabular}
\caption{Production of energy}
\label{tab_summary_energy_prod}
\end{table}

With 
\begin{equation*}
    \gamma=\frac{X_R+\beta\delta X_G}{X_R}-1=\beta\delta\frac{X_G}{X_R}.
\end{equation*}

Throughout the appendix, we assume gas and renewable prices remain constant
\begin{equation*}
    \hat{p}_G=\hat{p}_R=0.
\end{equation*}


The prices of the final, useful, and primary energy from the gas sector and renewables sector are noted
\[p_G^f, p_R^f,p_G^u, p_R^u,p_G^p, p_R^p. \]
The cost of final electricity is the same whether it comes from renewables or gas, 
\begin{equation}
\label{eq:price_elec}
\begin{cases}
     p_G^f&=p_R^f=p_R^p\\
      p_G^u&=p_R^u=\frac{p_R^f}{\alpha}\\
    p_G^f&=\frac{p_G^p}{\beta}
\end{cases}
\end{equation}


\subsection{Results summary}

In this paper, we focus on final energy, but within a framework similar to that of primary energy, where the energy sector delivers both electricity and fossil fuels (gas/oil). Note that in this appendix, ``final energy'' refers to final electricity. Electrification of energy productions can have diverse impacts on productivity at the sector level depending whether we are in a national account framework or in our $\zeta$ framework (\ER{Nom non contractuel}) developed for this paper, in a two or three-sector economy, or considering primary, final or useful energy.

The first result is that when considering final or useful energy, electrification does not alter the aggregate productivity in all cases (Table \ref{table_summary_results_energy}). 

The only problematic case arises when primary energy is analysed within a broad energy sector that includes both electricity and fossil fuels in our $\zeta$ framework. In this case, studying the energy sector alone correctly captures the sector's productivity growth. However, the aggregate productivity growth is non-zero and does not align with any of the others cases. 

The conclusion of Appendix (\er{A7}) where we consider productivity changes from an upstream sector does not hold. In Appendix (\er{A7}), the relationship $\hat{X}_M = \frac{\lambda_U}{\lambda_M} \hat{X}_U$ ensures that changes in the energy sector’s output are reflected in the downstream manufacturing sector. This mechanism does not hold for the energy sector itself, because useful energy remains unchanged—making it the only sector where a reduction in input can still provide exactly the same service.

\begin{table}[H]
\centering
\begin{tabular}{llll|lll}
\toprule
              & \multicolumn{3}{c}{National accounting} & \multicolumn{3}{c}{Framework} \\ \midrule
              & Primary      & Final      & Useful     
              & Primary   & Final   & Useful  \\ 
%\rowcolor[HTML]{D0CECE} 
\midrule
\textbf{\textit{Three sectors}} &              &            &             &           &         &         \\\midrule
Gas           & 0         & 0       &  0          
&   $\zeta \ln(1-\delta)$       &     $\zeta \ln(1-\delta)$      &    $\zeta \ln(1-\delta)$         \\
Renewables    & 0         & 0       &  0           
&      $\zeta \ln(\gamma)$       &       $\zeta \ln(\gamma)$     &   $\zeta \ln(\gamma)$        \\
Manufacturing & 0         & 0       &  0           
&     0     &  0        &     0    \\
Aggregate     & 0         & 0       &  0         
&    0      & 0         &   0      \\ 
%\rowcolor[HTML]{D0CECE} 
 \midrule
\textbf{\textit{Two sectors}}   &              &            &             &           &         &         \\\midrule
Energy        &    $\hat{A}_E=\hat{X}_E=-\hat{p}_E$    & 0       &  0          
&       $\hat{A}_E=\hat{X}_E=-\hat{p}_E$    &   0      &   0      \\
Manufacturing &  $\hat{A}_M=-a_{EM}\hat{X}_E$         & 0       &  0         
&      0     &    0     &       0  \\
Aggregate    & 0         & 0       &  0         
&     $\hat{A}=\frac{p_EX_E}{Y}\hat{X}_E$      &   0      &  0  \\
\bottomrule
\end{tabular}
\caption{Summary results}
\label{table_summary_results_energy}
\end{table}

Conclusion: If we adopt the final energy approach while including primary fossil fuels -— considering both the energy sector and the downstream sectors -— we encounter a quantification challenge. To ensure consistency, all energy flows would need to be converted into either final electricity or useful energy.\footnote{
\ER{François: Note that in the case displayed in the TPI version of the paper, the productivity growth we compute is the \emph{true} productivity growth of the energy sector computed with the national accounting primal/dual approach. But it would not be consistent to use it with other sectors.}
}

\textbf{The correct approach, when analysing a broad energy sector is to use the useful energy metric, ensuring that pure electrification does not affect productivity change in any sector.}

\ER{François: in other words, we will be fine with the useful energy approach within our framework.}


\subsection{Why is it different from previous appendix? (months ago)}
\ER{Subsection to be deleted after François' reading}

NB: in the previous appendix, why did I find it didn't change anything ? Because I was using our framework approach for the energy sector (with the subcase of a full switch, thus $\delta=1$) 
\[ \hat{A}_E=\hat{X}_E=\ln(\delta \beta )=\ln(\beta ),\]
However, for the manufacturing sector I was using the true productivity measure, 

  \begin{equation*}
    \left\{
    \begin{aligned}
        \hat{\Sigma}_X &= w_{EM}\hat{X}_{ME} 
                      = w_{EM}\ln \beta , \\
        \hat{\Sigma}_p &= w_{EM}\hat{p}_E
                      = w_{EM}\bigl(-\ln \beta \bigr).
    \end{aligned}
    \right.
\end{equation*}

where $w_{EM}$ is the energy share in M's cost.  Assuming $\hat{X}_M = 0$, one derives $\hat{p}_M = 0$ and
\begin{equation*}
    \hat{A}_M = -w_{EC}\ln \beta .
\end{equation*}
And then, in the aggregate, it cancels out. 
\[\hat{A}=0\]

But if we use the framework approach, then because $\hat{C}_M=0$\footnote{Because$w_{EC}$  remains constant, as the manufacturing sector's total energy bill does not change: 
\begin{equation*}
     p_GX_G 
         = \frac{p_R}{\beta }\beta X_G 
         = p_RX_R.
\end{equation*} Hence E's share in M's costs stays the same before and after.
}

\[
\left\{
\begin{aligned}
\hat{X}_M&=0\\
\hat{p}_M &= \hat{C}_M-\hat{X}_M= 0\\
\hat{A}&=0\\
\end{aligned}
\right.
\]

\newpage





\section{Three sectors: National Accounting and Our Framework}


In the following section we develop the computations of the results available in Table \ref{table_summary_results_energy}. 

We consider 3 sectors: Gas, Renewables and Manufacturing. In the next section we consider only 2: Energy and Manufacturing.

\subsection{Case 1: Primary-Energy Approach}
\label{subappendix_primary_energy}
\subsubsection{National Accounting Derivation}

In this approach, energy is measured in primary form. After a partial switch, the manufacturing sector M decreases its gas demand from $X_G$ to $(1-\delta)X_G$, while its use of renewables increases from $X_R$ to $X_R + \delta \beta X_G = X_R(1+\gamma)$.

\paragraph{Gas industry G}
Gas output shrinks by the factor $(1-\delta)$:
\[
    \hat{X}_G^p
= \ln\bigl((1-\delta) X_G\bigr)  -  \ln(X_G)
= \ln(1-\delta).
\]
From Eq.~\ref{eq:primal_appendix}, we assume all inputs to G scale down proportionally, so each input grows by $\ln(1-\delta)$. Using the same argument in the dual approach (Eq.~\ref{eq:dual_appendix}) would give $\hat{p}_G=0$, hence
\[
    \hat{A}_G
= \ln(1-\delta)
  - \sum_{j=1}^N a_{jG}\ln(1-\delta)
  - \sum_{f=1}^F \tilde{l}_{fG}\ln(1-\delta)
= \ln(1-\delta)  -  \ln(1-\delta)
= 0.
\]

\paragraph{Renewables industry R}
The renewables industry's output of final electricity rises from $X_R$ to $X_R(1+\gamma)$, so
\[
    \hat{X}_R^p
= \ln(1+\gamma).
\]
Similarly, scaling all inputs proportionally yields
\[
    \hat{A}_R = 0.
\]

\paragraph{Manufacturing sector M}
Its gas input growth is $\hat{X}_G^p=\ln(1-\delta)$, and its renewables input growth is $\hat{X}_R^p=\ln(1+\gamma)$. With $\hat{X}_M=0$ and no change in labour usage, the primal approach (Eq.~\ref{eq:primal_appendix}) gives
\[
    \hat{A}_M
= 0
  - \bigl[a_{GM}\ln(1-\delta)
     + a_{RM}\ln(1+\gamma)\bigr]
  - 0,
\]
which is equal to 0. 

We verify this by looking at M's expenditures. Since M's final energy requirement and the price of electricity remain unchanged, the total energy cost before and after the switch satisfies \(C_M = C_M'\). Let \(C_G = p_G^pX_G\) and \(C_R = p_R^pX_R\), so \(C_M = C_G + C_R\). After the switch,
\[
   C_G' = C_G(1-\delta), \quad
   C_R' = C_R(1+\gamma),\]
   Hence
\begin{equation}
\label{eq_cost_argument_M}
\begin{aligned}
         C_M'&=C_G'+C_R'=C_G  (1-\delta)+C_R(1+\gamma)\\&
     =p_G^pX_G  (1-\delta)+p_R^pX_R\left(1+\beta\delta\frac{X_G}{X_R}\right)\\
    &= p_G^pX_G  (1-\delta)+p_R^pX_R+\frac{p_G^p}{\beta}\beta \delta X_G\\
     &= p_G^pX_G +p_R^pX_R
     =C_M.
\end{aligned}
\end{equation}

It follows that
\begin{equation}
 \label{eq_weighted_share_change_cost}
\begin{aligned}
a_{GM} \ln(1-\delta) + a_{RM} \ln(1+\gamma)&= 
a_{GM}\left[ \ln(C_G') - \ln(C_G) \right] + 
a_{RM} \left[ \ln(C_R') - \ln(C_R) \right]\\&= 
\frac{C_G}{C_G + C_R} \hat{C}_G' + 
\frac{C_R}{C_G + C_R} \hat{C}_R'\\
  &=\widehat{(C_G+C_R)}\\
 &=   \ln(C_G' + C_R') - \ln(C_G + C_R) \\
&=\ln(T') - \ln(T)  \\
&= 0.
\end{aligned}
\end{equation}
because the total cost $C_M$ before and after the switch remains the same,\footnote{
Noting that the log growth of $x+y$, $\widehat{x+y}$, is $\frac{x}{x+y} \hat{x} + \frac{y}{x+y} \hat{y}$.
}
 hence
\[
    \hat{A}_M=0.
\]

\paragraph{Aggregate}
With Domar weights $\lambda_G,\lambda_R,\lambda_M$, aggregate productivity growth is
\[
    \hat{A}
= \lambda_G\hat{A}_G
  + \lambda_R\hat{A}_R
  + \lambda_M\hat{A}_M
= 0.
\]

\subsubsection{Framework Derivation}

Here, we compute industry-level productivity growth using
\[
    \hat{A}_i
= \zeta \hat{X}_i  +  (\zeta-1) \hat{p}_i.
\]

\paragraph{Gas and Renewables}
Since gas output shrinks by $\ln(1-\delta)$ while $\hat{p}_G=0$, we get
\[
    \hat{X}_G^p = \ln(1-\delta),
    \quad
    \hat{A}_G = \zeta  \hat{X}_G^p = \zeta \ln(1-\delta).
\]
Similarly, renewables output rises by $\ln(1+\gamma)$ with $\hat{p}_R=0$, so
\[
    \hat{A}_R
= \zeta \ln(1+\gamma).
\]

\paragraph{Manufacturing M}
Manufacturing output does not grow ($\hat{X}_M=0$). From the cost argument (Eq. \ref{eq_cost_argument_M}), the price $\hat{p}_M$ also stays at zero. Thus
\[
    \hat{A}_M = 0.
\]

\paragraph{Aggregate}
Finally, the Domar-weighted sum gives
\[
    \hat{A}
= \lambda_G \zeta \ln(1-\delta)
   + 
  \lambda_R \zeta \ln(1+\gamma)
= 0.
\]

\subsection{Case 2: Final-Electricity Approach}
\label{subappendix_final_energy}
\subsubsection{National Accounting Derivation}

After the partial switch, M's demand for gas-based electricity decreases by the factor $(1-\delta)$, and its demand for renewables rises by a factor $(1+\gamma)$, with exactly the same substitution condition $\beta \delta X_G = X_R \gamma$. All inputs of both G and R are scaled proportionally to their respective outputs.

\paragraph{Gas industry G}
\[
    \hat{X}_G^f
= \ln\bigl(\beta(1-\delta)X_G\bigr)
  - \ln(\beta X_G)
= \ln(1-\delta),
  \quad
  \hat{A}_G=0.
\]

\paragraph{Renewables industry R}
\[
    \hat{X}_R^f = \ln(1+\gamma),
    \quad
    \hat{A}_R = 0.
\]

\paragraph{Manufacturing M}
The price of electricity remains unchanged regardless of generation method, so substituting renewables for gas does not alter cost shares or labour usage.\footnote{For the sake of completeness, one can explicitly derive this result using the same method as in the previous case. With $\hat{X}_M=0$ and no change in M's labour usage} Hence
\[
    \hat{A}_M=0.
\]

\paragraph{Aggregate}
Again,
\begin{equation}
      \hat{A}
= \lambda_G\hat{A}_G
  + \lambda_R\hat{A}_R
  + \lambda_M\hat{A}_M
=0.
\end{equation}
\subsubsection{Framework Derivation}

Since $\hat{X}_G^f=\ln(1-\delta)$, $\hat{p}_G=0$, $\hat{X}_R^f=\ln(1+\gamma)$, $\hat{p}_R=0$, we obtain
\[
    \hat{A}_G = \zeta \ln(1-\delta),
    \quad
    \hat{A}_R = \zeta \ln(1+\gamma).
\]
Because the electricity price does not change, in manufacturing $\hat{X}_M=0$ and $\hat{p}_M=0$, so
\[
    \hat{A}_M=0,
\]
Hence, aggregate productivity growth is
\begin{equation}
    \hat{A}= \lambda_G \zeta \ln(1-\delta) + \lambda_R \zeta \ln(1+\gamma)
=0.
\end{equation}

\subsection{Case 3: Useful-Energy Approach}
\label{subappendix_useful_energy}
\subsubsection{National Accounting Derivation}

Now G and R directly sell \emph{useful} electricity to M. The partial switch replaces $\delta \alpha\beta X_G$ gas-based electricity by $\delta \alpha\beta X_G$ renewables-based electricity, leaving total useful energy demand unchanged. Both industries scale their inputs proportionally, so
\[
    \hat{X}_G^u
= \ln(1-\delta),
    \quad
    \hat{X}_R^u
= \ln(1+\gamma),
    \quad
    \hat{A}_G=0,
    \quad
    \hat{A}_R=0.
\]
In manufacturing, $\hat{X}_M=0$ and labour is unchanged, so
\[
   \hat{A}_M
= - \bigl[a_{GM}\ln(1-\delta)
       + a_{RM}\ln(1+\gamma)\bigr]
=0.
\]
Hence, 
\[
   \hat{A}=0.
\]

\subsubsection{Framework Derivation}

Exactly the same substitution argument applies; we find
\[
    \hat{X}_G^u
= \ln(1-\delta), \hat{p}_G=0,
    \quad
    \hat{X}_R^u
= \ln(1+\gamma), \hat{p}_R=0,
\]
and so
\[
    \hat{A}_G=\zeta \ln(1-\delta),
    \quad
    \hat{A}_R=\zeta \ln(1+\gamma).
\]
Since M's total energy and price do not change,
\[
    \hat{X}_M=0,\quad\hat{p}_M=0,\quad \hat{A}_M=0.
\]
Domar-aggregating again gives $\hat{A}=0$.

\section{Two sectors: National Accounting and Our Framework}


In this section we develop the computations of the results available in Table \ref{table_summary_results_energy} with only 2 sectors: Energy and Manufacturing.

\subsection{Case 4: Primary-Energy with Two Industries}
\label{subappendix_primary_energy_two_sectors}
\subsubsection{National Accounting Derivation}

Here, an ``energy industry E'' delivers two types of primary outputs to manufacturing, gas-based and renewables-based, whose total final use remains constant in terms of useful output, but not necessarily in the exact energy content. Denote $X_E^p=X_G^p + X_R^p$ in primary terms and assume $p_R^p = p_G^p/\beta$. The partial switch reduces $X_G^p$ from $X_G$ to $(1-\delta)X_G$ and increases $X_R^p$ from $X_R$ to $X_R + \beta\delta X_G$, which keeping total energy sales $C_E$ constant. 
We find
\begin{equation}
\label{eq_constant_CE}
    \begin{aligned}
    \hat{C}_E&=\ln(p_G(1-\delta)X_G+p_R(X_R+X_G\delta \beta ))-\ln(p_GX_G+p_RX_R)\\
    &=\ln(p_GX_G+p_G\frac{X_R}{\beta })-\ln(p_GX_G+p_G\frac{X_R}{\beta })\\
&=0.
\end{aligned}
\end{equation}
Hence,
\begin{equation}
\label{eq_pe_xe}
     \hat{C}_E
= \hat{X}_E^p + \hat{p}_E^p
= 0
\quad\Longrightarrow\quad
    \hat{p}_E^p
= - \hat{X}_E^p.
\end{equation}

From the primal and dual definition of productivity growth
\begin{equation}
    \hat{A}_E = - \hat{p}_E^p =\hat{X}_E^p.
\end{equation}


In manufacturing, $p_EX_E^p + \sum_f w_fL_f = p_M X_M$. Since the cost of energy is constant (Eq. \ref{eq_constant_CE}), the cost share of energy in M's expenses is constant, so
\begin{equation*}
        \hat{p}_M=\hat{C}_M-\hat{X}_M
= 0
\end{equation*}

The manufacturing sector uses energy as an intermediate input and factors. 
From the primal definition of productivity, Eq. \ref{eq:primal_appendix},\footnote{The dual definition Eq. \ref{eq:dual_appendix} gives the same result.}
\begin{equation*}
   \hat{A}_M= \hat{X}_M-a_{EM}\hat{X}_E^p= -a_{EM}\hat{X}_E^p.
\end{equation*}

Finally, the Domar-weighted sum over E and M cancels out to give
\begin{equation}
    \begin{aligned}
    \hat{A}&=\lambda_E \hat{A}_E+\lambda_M \hat{A}_M\\
    &=\frac{X_Ep_E}{Y} \hat{X}_E^p+ \frac{p_XX_M}{Y} (-a_{EM}\hat{X}_E^p)\\
    &=\frac{X_Ep_E}{Y} \hat{X}_E^p - \frac{p_XX_M}{Y} \frac{X_Ep_E}{p_XX_M}\hat{X}_E^p\\
     &=\frac{X_Ep_E}{Y} \hat{X}_E^p - \frac{X_Ep_E}{Y}\hat{X}_E^p\\
     &=0.
\end{aligned}
\end{equation}

\subsubsection{Framework Derivation}

Since, $\hat{p}_E= - \hat{X}_E^p.$ (Eq. \ref{eq_pe_xe}),
\begin{equation*}
      \hat{A}_E
= \zeta \hat{X}_E^p + (\zeta-1) \hat{p}_E
= \zeta \hat{X}_E^p + (\zeta-1)\bigl(-\hat{X}_E^p\bigr)
= \hat{X}_E^p.
\end{equation*}

In manufacturing, $\hat{X}_M=0$, $\hat{p}_M=0$, thus
\[
    \hat{A}_M=0.
\]
Hence the Domar-weighted sum is
\[
    \hat{A}
= \lambda_E \hat{A}_E
  + \lambda_M \hat{A}_M
= \lambda_E \hat{X}_E^p
\]


\subsection{Case 5: Final- or Useful-Energy Approach with Two Industries}
\label{subappendix_final_energy_two_sectors}
\subsubsection{National Accounting Derivation}

Now E delivers \emph{final} (or \emph{useful}) electricity from both gas and renewables, and M's demand for total final (or useful) energy is constant across the partial switch. The price of final energy is total expenditure on energy divided by total units of energy. Since (i) total energy use is unchanged and (ii) total expenditure on energy is also unchanged (the decline in gas spending exactly matches the rise in renewables spending), the overall price of energy must stay constant. 
Therefore,
\[
    \hat{X}_E^f
= 0,
    \quad
    \hat{p}_E= 0,
\]
hence
\[
    \hat{A}_E=0.
\]
Since the manufacturing sector faces the same final-energy demand at the same energy price,
\[
    \hat{A}_M=0,
\]
and thus $\hat{A}=0$.

\subsubsection{Framework Derivation}

If the energy industry changes its mix of inputs but the total final (or useful) output is unchanged, then
\[
   \hat{X}_E^f = 0,
   \quad
   \hat{p}_E = 0
    \Longrightarrow 
   \hat{A}_E=0.
\]
In manufacturing, neither quantity nor price changes, so $\hat{X}_M=0$ and $\hat{p}_M=0$, yielding $\hat{A}_M=0$. Domar-weighting gives $\hat{A}=0$ overall.



\subsection{Price dynamics - Two sectors, primary energy}

Now we consider that the price of gas and renewable are not fixed.

\bigskip

\subsubsection{ National Accounting Approach}

\paragraph{Energy Sector \((E)\)}

\noindent
\textbf{General formula for the log-growth of }\(p_E\):
\[
  p_E 
   =  
  \frac{p_G X_G  +  p_R X_R}{ X_E },
  \quad
  \hat{p}_E  =  \ln\!\biggl(\frac{p_E(t+1)}{ p_E(t)\!}\biggr).
\]
By taking the logarithmic derivative of the numerator minus that of the denominator, one obtains:
\[
  \hat{p}_E
   = 
  \underbrace{
    \widehat{\bigl(p_G X_G + p_R X_R\bigr)}
  }_{\text{log-growth of the sum}}
   - 
  \underbrace{
    \hat{X}_E
  }_{\text{log-growth of output}}
   .
\]
We get
\[
    \hat{p}_E
     = 
    \frac{p_G X_G}{p_E X_E} 
      (\hat{p}_G + \hat{X}_G)
     + 
    \frac{p_R X_R}{p_E X_E} 
      (\hat{p}_R + \hat{X}_R)
     - 
    \hat{X}_E.
\]

From the dual,
\[
  \hat{A}_E
   = 
  - \Bigl[
        \hat{p}_E
         \bigr)
      \Bigr],
\]
Because we only consider one sector: E. And the price of the factors don't change.

\begin{align*}
     \hat{A}_E
   &= 
  - \Bigl[
        \frac{p_G X_G}{p_E X_E} 
      (\hat{p}_G + \hat{X}_G)
     + 
    \frac{p_R X_R}{p_E X_E} 
      (\hat{p}_R + \hat{X}_R)
     - 
    \hat{X}_E
\end{align*}


\paragraph{Manufacturing Sector \((M)\)}

Manufacturing $M$ buys $X_E$ at price $p_E$, plus factors (labour, capital) whose usage and prices do not change. Let $a_{EM}$ be the cost share of energy. Then:

\paragraph{Primal.}
\[
  \hat{A}_M
   = 
  \hat{X}_M
   - 
  \bigl[a_{EM} \hat{X}_E
     + 
    \sum_{f}\tilde{l}_{fM} \hat{L}_{fM}\bigr].
\]
Since the factors do not change ($\hat{L}_{fM}=0$), this reduces to
\[
  \hat{A}_M
   = 
  \hat{X}_M 
   - 
  a_{EM} \hat{X}_E.
\]
If $M$'s output also stays unchanged ($\hat{X}_M=0$), then $\hat{A}_M=- a_{EM} \hat{X}_E$. 

From the dual,
\[
  \hat{A}_M
   = 
  - \Bigl[ 
      \hat{p}_M
       - 
      a_{EM} \hat{p}_E
    \Bigr],
\]
again ignoring factor-price terms that are zero when $\hat{w}_f=0$.

\begin{equation*}
      \hat{p}_M=\hat{C_M}=\frac{p_EX_E}{p_MX_M}\hat{C_E}=\frac{p_EX_E}{p_MX_M}\left(\hat{X}_E+\hat{p}_E\right)
\end{equation*}

\paragraph{Aggregate}

Let $\lambda_E$ and $\lambda_M$ be the Domar weights. Then
\[
  \hat{A}
   = 
  \lambda_E \hat{A}_E
   + 
  \lambda_M \hat{A}_M.
\]
Explicitly, using dual forms:
\begin{align*}
     \hat{A}
   &= 
  - \Bigl[
    \lambda_E \bigl(\hat{p}_E\bigr)
     + 
    \lambda_M \bigl(\hat{p}_M - a_{EM} \hat{p}_E\bigr)
  \Bigr]\\
 &=  - \Bigl[\frac{p_EX_E}{Y} \hat{p}_E
     + 
   \frac{p_MX_M}{Y}\bigl(\hat{p}_M - \frac{p_EX_E}{p_MX_M} \hat{p}_E\bigr)
  \Bigr]\\
  &=  - \Bigl[\frac{p_EX_E}{Y} \hat{p}_E
     + 
   \frac{p_MX_M}{Y}\bigl(\frac{p_EX_E}{p_MX_M}\left(\hat{X}_E+\hat{p}_E\right) - \frac{p_EX_E}{p_MX_M} \hat{p}_E\bigr)
  \Bigr]\\
  &= - \frac{p_EX_E}{Y}\Bigl[\hat{p}_E 
  + \left(\hat{X}_E+\hat{p}_E\right) - \hat{p}_E
  \Bigr]\\
    &= - \frac{p_EX_E}{Y}\Bigl[
  \left(\hat{X}_E+\hat{p}_E\right)
  \Bigr]\\
  &=
  - \frac{p_EX_E}{Y}\Bigl[
  \left(\hat{X}_E+\frac{p_G X_G}{p_E X_E} 
      (\hat{p}_G + \hat{X}_G)
     + 
    \frac{p_R X_R}{p_E X_E} 
      (\hat{p}_R + \hat{X}_R)
     - 
    \hat{X}_E.\right)
  \Bigr]\\
   &=
  - \frac{p_EX_E}{Y}\Bigl[
  \left(\frac{p_G X_G}{p_E X_E} 
      (\hat{p}_G + \hat{X}_G)
     + 
    \frac{p_R X_R}{p_E X_E} 
      (\hat{p}_R + \hat{X}_R)
    \right)
  \Bigr]\\
\end{align*}

Only the decrease in the price of energy appear, as if the electrification didn't take place (0 productivity gain).


\subsubsection{Framework Approach}

In this alternative definition, each sector $i$ has TFP growth
\[
  \hat{A}_i
   = 
  \zeta \hat{X}_i
   + 
  (\zeta - 1) \hat{p}_i,
\]
where $\zeta$ is a parameter (e.g.\ describing returns to scale).

\paragraph{Energy Sector \((E)\)}

\[
  \hat{A}_E
   = 
  \zeta \hat{X}_E 
   + 
  (\zeta - 1) \hat{p}_E.
\]

\paragraph{Manufacturing Sector \((M)\)}

\[
  \hat{A}_M
   = 
  \zeta \hat{X}_M
   + 
  (\zeta - 1) \hat{p}_M.
\]
With $\hat{X}_M=0$ 

\paragraph{Aggregate}

Again, Domar weighting:
\[
  \hat{A}
   = 
  \lambda_E \hat{A}_E
   + 
  \lambda_M \hat{A}_M,
\]
so



\begin{align*}
    \hat{A}
   &= 
  \lambda_E\bigl[\zeta \hat{X}_E + (\zeta -1) \hat{p}_E\bigr]
   + 
  \lambda_M\bigl[(\zeta -1) \hat{p}_M\bigr]\\
   &= 
  \lambda_E\bigl[\zeta \hat{X}_E + (\zeta -1) \hat{p}_E\bigr]
   + 
  \lambda_M\bigl[(\zeta -1) \frac{p_EX_E}{p_MX_M}\left(\hat{X}_E+\hat{p}_E\right)\bigr]\\
     &= 
  \lambda_E\bigl[\zeta \hat{X}_E + (\zeta -1) \hat{p}_E
   + 
(\zeta -1) \left(\hat{X}_E+\hat{p}_E\right)\bigr]\\
     &= 
  \lambda_E\bigl[(2\zeta-1) \hat{X}_E + 2(\zeta -1) \hat{p}_E \bigr]\\
\end{align*}

We get the same result as in the previous section if $\hat{X}_E=-\hat{p}_E $.

\end{document}


\newpage

\section{National accounting framework}
\subsection{Case 1: Primary-Energy Approach}
\label{subappendix_primary_energy}

This appendix provides the detailed derivation of industry-level and aggregate productivity changes when we measure energy flows in primary energy rather than final or useful energy. 

After the partial switch, M's demand for gas decreases from $X_G$ to $(1-\delta) X_G$, while its use of renewable electricity increases from $ X_R $ to $X_R + \delta \beta X_G=X_R(1+\gamma) $.


\subsubsection{Gas industry G} The gas industry is delivering gas.
Its output shrinks by the factor $(1-\delta)$.  Thus
\begin{equation*}
    \hat{X}_G^p
= \ln\left((1-\delta) X_G\right) - \ln(X_G)
= \ln(1-\delta).
\end{equation*}
Applying Eq. \ref{eq:primal_appendix}, we assuming all inputs to G are scaled down proportionally with output  $\hat{X}_G^p = \ln(1-\delta)$, so each $\hat{X}_{jG}= \ln(1-\delta)$ and $\hat{L}_{fG}= \ln(1-\delta)$. Thus\footnote{
One can check that using the dual approach from Eq. \ref{eq:dual_appendix} gives the same results as $\hat{p}_G=0$, as well as $\forall j,f~\hat{p}_j=\hat{w}_f=0$.
}
\begin{equation}
    \hat{A}_G = \ln(1-\delta) -  \sum_{j=1}^N a_{ji} \ln(1-\delta)  -  \sum_{f=1}^F \tilde{l}_{fi} \ln(1-\delta)=  \ln(1-\delta) - \ln(1-\delta)=0
\end{equation}

\subsubsection{Renewables industry R}
The renewables industry's output of final electricity rises from $X_R$ to $X_R(1+\gamma)$. 
Hence
\begin{equation*}
    \hat{X}_R^p = \ln(1+\gamma),
\end{equation*}
and by the same reasoning as above, scaling all inputs proportionality with output,
\begin{equation}
    \hat{A}_R = 0.
\end{equation}

\subsubsection{Manufacturing M} The manufacturing sector undergoes a change in input, 
\begin{equation*}
    \sum_{j\in\{G,R\}} a_{jM} \hat{X}_{jM}
= a_{GM} \hat{X}_G^p + a_{RM} \hat{X}_R^p,
\end{equation*}
where
\begin{equation*}
    \hat{X}_G^p= \ln(1-\delta),
\quad \text{and}\quad 
\hat{X}_R^p= \ln(1+\gamma).
\end{equation*}

With $\hat{X}_M=0$ and no change in M's labour usage, Eq. \ref{eq:primal_appendix} implies
\begin{equation}
\label{eq_final_Ahat_M}
    \hat{A}_M
= \hat{X}_M 
  - \sum_{j\in\{G,R\}} a_{jM} \hat{X}_{jM}
  - \sum_f \tilde{l}_{fM} \hat{L}_{fM}
= 0
  - \bigl[a_{GM} \ln(1-\delta) + a_{RM} \ln(1+\gamma)\bigr]
  - 0.
\end{equation}
which is equal to 0. 

\vspace{3mm}
We derive this result by explicitly considering M's expenditures. Since the total quantity of final energy required by M remains constant (Eq. \ref{eq_constant_final_demand_elec}), and the price of electricity does not change (Eq. \ref{eq:price_elec}), the total cost of manufacturing remains unchanged before and after the switch, i.e., \( C_M = C_M' \).   Let 
$C_G = p_G^p  X_G$, 
$C_R = p_R^p  X_R$,
and
$C_M = C_G + C_R$
denote the initial monetary expenditures on gas, renewable electricity, and total energy, respectively. After the switch, these expenditures become  $C_G'=C_G  (1-\delta),~C_R'= C_R(1+\gamma)$, and 

\begin{align*}
    C_M'&=C_G'+C_R'=C_G  (1-\delta)+C_R(1+\gamma)\\
   
\end{align*} 

We then write
\footnote{
Noting that the log growth of $x+y$, $\widehat{x+y}$, is $\frac{x}{x+y} \hat{x} + \frac{y}{x+y} \hat{y}$.
}


\begin{equation}
 \label{eq_weighted_share_change_cost}
\begin{aligned}
a_{GM} \ln(1-\delta) + a_{RM} \ln(1+\gamma)&= 
a_{GM}\left[ \ln(C_G') - \ln(C_G) \right] + 
a_{RM} \left[ \ln(C_R') - \ln(C_R) \right]\\&= 
\frac{C_G}{C_G + C_R} \hat{C}_G' + 
\frac{C_R}{C_G + C_R} \hat{C}_R'\\
  &=\widehat{(C_G+C_R)}\\
 &=   \ln(C_G' + C_R') - \ln(C_G + C_R) \\
&=\ln(T') - \ln(T)  \\
&= 0.
\end{aligned}
\end{equation}

Hence, 
\begin{equation}
    \hat{A}_M=0
\end{equation}


\subsubsection{Aggregate TFP}
Let $\lambda_G,\lambda_R,\lambda_M$ be the Domar weights.  Aggregate productivity growth $\hat{A}$ is
\begin{equation}
\label{eq_final_aggregate_tfp}
    \hat{A} \equiv
 \lambda_G \hat{A}_G
  + \lambda_R \hat{A}_R
  + \lambda_M \hat{A}_M=0.
\end{equation}







\subsection{Case 2: Final-Electricity Approach}
\label{subappendix_final_energy}

After the partial switch, M's demand for gas-based electricity decreases from $\beta X_G $ to $\beta(1-\delta) X_G $, while its use of renewable electricity increases from $ X_R $ to $ X_R (1+\gamma) $. Note that the switched quantity of gas-based electricity, is exactly equal to the additional amount of renewable electricity
\begin{equation}
    X_R\gamma=X_R \beta\delta\frac{X_G}{X_R}=\beta \delta X_G.
\end{equation}

\subsubsection{Gas Industry G}

The gas industry G shrink its output by a factor  $(1-\delta)$. Thus
\begin{equation*}
    \hat{X}_G^f 
= \ln\left((1-\delta)\beta X_G\right) - \ln(\beta X_G)
= \ln(1-\delta).
\end{equation*}

Applying Eq. \eqref{eq:primal_appendix}, we assuming all inputs to G are scaled down proportionally with output, thus
\begin{equation}
    \hat{A}_G = 0.
\end{equation}


\subsubsection{Renewables Industry R}

R's output rises by a factor $\gamma$.  
Hence
\[
\hat{X}_R^f= \ln(1+\gamma).
\]

Similarly, with all input usage scaling proportionally,
\[
\hat{A}_R = 0.
\]




\subsubsection{Manufacturing M}
There are two ways to analyse the manufacturing sector.
Since the price of electricity remains the same regardless of whether it is generated from gas or renewables (Eq. \ref{eq:price_elec}), substituting one for the other does not alter cost shares, input use, the final product’s price, or production levels. Consequently, the impact on productivity is null.

For the sake of completeness, one can explicitly derive this result using the same method as in the previous case. With $\hat{X}_M=0$ and no change in M's labour usage, Eq. \ref{eq:primal_appendix} implies
\begin{equation}
\label{eq_final_Ahat_M}
    \hat{A}_M
= \hat{X}_M 
  - \sum_{j\in\{G,R\}} a_{jM} \hat{X}_{jM}
  - \sum_f \tilde{l}_{fM} \hat{L}_{fM}
= 0
  - \bigl[a_{GM} \ln(1-\delta) + a_{RM} \ln(1+\gamma)\bigr]
  - 0=0.
\end{equation}


\subsubsection{Aggregate TFP}
Let $\lambda_G,\lambda_R,\lambda_M$ be the Domar weights. Thus, aggregate productivity growth  $\hat{A}$ is
\begin{equation*}
    \hat{A} \equiv
 \lambda_G \hat{A}_G
  + \lambda_R \hat{A}_R
  + \lambda_M \hat{A}_M=0.
\end{equation*}



\subsection{Case 3: Useful-Energy Approach}
\label{subappendix_useful_energy}

The Gas and Renewables industry now directly sell useful electricity from gas and renewables respectively. The useful energy demand by M does not change with the switch from gas to renewable.

The manufacturing sector initially purchases $\alpha\beta X_G$ units of gas-based useful electricity and $\alpha X_R$ units of renewables-based useful electricity.  We consider a small partial switch: a fraction $\delta \alpha\beta X_G$ of M's gas-based useful electricity is replaced by renewables. 

As in the previous case, the inputs of both the gas and renewables sectors are scaled up or down proportionally with output
\begin{equation}
    \hat{X}_G^u=\ln(1-\delta),\quad\text{and}\quad \hat{X}_R^u=\ln(1+\gamma).
\end{equation}
Thus
\begin{equation}
    \hat{A}_G=0,\quad\text{and}\quad \hat{A}_R=0.
\end{equation}


\subsubsection{Manufacturing M}

Since $\hat{X}_M=0$ and we assume no change in M's labour usage,
\[
\hat{A}_M
= \hat{X}_M
  - \sum_{j\in\{G,R\}} a_{jM} \hat{X}_{jM}
  - \sum_f \tilde{l}_{fM} \hat{L}_{fM}
= 
  - \bigl[a_{GM} \ln(1-\delta) + a_{RM} \ln(1+\gamma)\bigr]=0.
\]

\subsubsection{Aggregate TFP}
Aggregate productivity growth is
\begin{equation}
    \hat{A}=0
\end{equation}


\subsection{Case 4: Primary with two sectors}

We consider an energy sector that supplies both gas and renewable electricity. The manufacturing sector consumes gas to generate its own electricity while also using electricity directly.

\subsubsection{Energy industry E} The energy sector generates gas-based electricity, which declines by a factor of $(1-\delta)$, while producing renewable electricity, which increases by a factor of $(1+\gamma)$. As a result, the total quantity of useful energy delivered to the manufacturing sector remains unchanged. To ease the reading, we then still refer to the production of each technology as $X_G$ and $X_R$, such as $X_E=X_R+X_G$.
Remember that,
\begin{equation*}
    p_R^f=p_G^f=p_R^p=\frac{p_G^p}{\beta}.
\end{equation*}
Thus, we compute, the energy price at time $t$,
\begin{align*}
 \ln(p_{E,t})=& -\left(\ln(X_G p_G+X_R p_R) - \ln(X_G+X_R)\right)\\
     &= -\ln(p_G)-\left(\ln(X_G +\frac{X_R}{\beta} ) - \ln(X_G+X_R)\right),
\end{align*}

\bigskip

\noindent the energy price at time $t+1$,
\begin{align*}
 \ln(p_{E,t+1}^p)&=
    \left(\ln\left[(1-\delta)X_G p_G+(X_R+\delta X_G \beta) \frac{p_G}{\beta}\right] - \ln((1-\delta)X_G+(X_R+X_G \delta \beta) )\right)\\
    &=
    \left(\ln\left[(1-\delta)X_G p_G+X_R\frac{p_G}{\beta}+\delta X_Gp_G\right] - \ln((1-\delta)X_G+X_R+X_G \beta\delta)\right)\\
    &=
   \ln(p_G) +\ln\left[(1-\delta)X_G+\frac{X_R}{\beta}+\delta X_G\right] - \ln((1-\delta)X_G+X_R+\delta X_G \beta)\\
    &=
   \ln(p_G) +\ln\left[X_G+\frac{X_R}{\beta}\right] - \ln((1-\delta)X_G+X_R+\delta X_G \beta),
\end{align*}

\bigskip

\noindent the logarithmic growth of the price of energy,
\begin{align*}
\hat{p}_E^p&=  \ln(p_G) +\ln\left[X_G+\frac{X_R}{\beta}\right] - \ln((1-\delta)X_G+X_R+\delta X_G \beta)\\
&\quad -\ln(p_G)-\left(\ln(X_G +\frac{X_R}{\beta} ) - \ln(X_G+X_R)\right)\\
&=   - \ln((1-\delta)X_G+X_R+\delta X_G \beta)+\ln(X_G+X_R).
\end{align*}

\bigskip

\noindent and the logarithmic growth of the production of energy,
\begin{align*}
    \hat{X}_E^p &=\ln((1-\delta) X_G + (X_R+\beta\delta X_G)-\ln(X_G +X_R)\\
    &=-\hat{p}_E^p.
\end{align*}

Which mean that total sales $C_E$ are constant, since
\begin{equation*}
    \hat{C}_E=\hat{p}_E^p+\hat{X}_E^p=0.
\end{equation*}

From the primal and dual definition of productivity growth
\begin{equation}
    \hat{A}_E = - \hat{p}_E^p =\hat{X}_E^p.
\end{equation}


\subsubsection{Manufacturing M} The manufacturing sector uses energy as an intermediate input and factors. 
From the primal definition of productivity, Eq. \ref{eq:primal_appendix},
\begin{equation*}
   \hat{A}_M= \hat{X}_M-a_{EM}\hat{X}_E^p= -a_{EM}\hat{X}_E^p.
\end{equation*}
Let us check that the cost share of energy are constant, on one hand we have,
\begin{align*}
\begin{cases}
    \text{Sales at time }t,\quad& X_Gp_G+X_Rp_R=p_G\left(X_G+\frac{X_R}{\beta}\right)\\
    \text{Sales at time }t+1,\quad&  (1-\delta)X_G p_G+(X_R+\beta\delta X_G)\frac{p_G}{\beta}=p_G\left( X_G+\frac{X_R}{\beta}\right),
\end{cases}
\end{align*}
on the other hand we have, since $\widehat{L_fw_f}=0$,
\begin{equation*}
    \hat{p}_M=\widehat{\left(\frac{p_EX_E+L_fw_f}{X_M}\right)}=\widehat{\left(p_EX_E+L_fw_f\right)}=
    \frac{p_EX_E}{p_EX_E+L_fw_f}\widehat{p_EX_E}=
    \frac{p_EX_E}{p_XX_M}(\hat{p}_E^p+\hat{X}_E^p)=0.
\end{equation*}
Therefore, the cost share of energy is constant over time.

Using the dual equation, Eq. \ref{eq:dual_appendix},
\begin{equation*}
    \hat{A}_M=-(\hat{p}_M-a_{EM}\hat{p}_E^p)=-a_{EM}\hat{X}_E^p.
\end{equation*}


\subsubsection{Aggregate} 
Aggregate productivity growth is
\begin{align*}
    \hat{A}&=\lambda_E \hat{A}_E+\lambda_M \hat{A}_M\\
    &=\frac{X_Ep_E}{Y} \hat{X}_E^p+ \frac{p_XX_M}{Y} (-a_{EM}\hat{X}_E^p)\\
    &=\frac{X_Ep_E}{Y} \hat{X}_E^p - \frac{p_XX_M}{Y} \frac{X_Ep_E}{p_XX_M}\hat{X}_E^p\\
     &=\frac{X_Ep_E}{Y} \hat{X}_E^p - \frac{X_Ep_E}{Y}\hat{X}_E^p\\
     &=0.
\end{align*}


\subsection{Case 5: Useful-Energy Approach for two industries}
\label{subappendix_final_energy}


\subsubsection{Energy industry} After and before the switch, the output from the energy industry remains the same in final energy, but the composition changes. 
\[
\hat{X}_E^u=0
\]

We compute the average price of energy before the switch ($t$) and after the switch ($t+1$), along with the logarithmic growth $\hat{p}_E$. 

\bigskip

\noindent\textbf{Energy price at time }$t$:
\begin{align*}
\ln(p_{E,t})
&= -\Bigl(\ln(\alpha\beta X_G p_G + \alpha X_R p_R)  -  \ln(\alpha\beta  X_G + \alpha X_R)\Bigr)\\
&= -\ln(p_G) - \Bigl(\ln\bigl(\alpha X_G + \tfrac{\alpha X_R}{\beta}\bigr)  -  \ln(\beta \alpha X_G +\alpha  X_R)\Bigr).
\end{align*}

\bigskip

\noindent\textbf{Energy price at time }$t+1$:
\begin{align*}
\ln(p_{E,t+1})
&= \ln\Bigl(\tfrac{X_{G,t+1} p_G + X_{R,t+1} p_R}{\beta X_{G,t+1} + X_{R,t+1}}\Bigr) \\
&= \ln\bigl(X_{G,t+1} p_G + X_{R,t+1} p_R\bigr)
    -  \ln\bigl( X_{G,t+1} + X_{R,t+1}\bigr)\\
&= \ln\Bigl[(1-\delta)\alpha \beta X_G p_G + \alpha X_R \tfrac{p_G}{\beta} + \delta \alpha X_G p_G\Bigr]
    -  \ln\Bigl[(1-\delta) \alpha\beta X_G + X_R + \delta\alpha \beta X_G\Bigr]\\
&= \ln(p_G)
    +  \ln\Bigl[(1-\delta)\alpha X_G + \tfrac{X_R}{\beta} + \delta X_G\Bigr]
    -  \ln\Bigl[(1-\delta) \beta X_G + X_R + \delta \beta X_G\Bigr]\\
&= \ln(p_G)
    +  \ln\Bigl[X_G + \tfrac{X_R}{\beta}\Bigr]
    -  \ln\bigl(\beta X_G + X_R\bigr).
\end{align*}

\bigskip

\noindent\textbf{Logarithmic growth of the price of energy}:
\begin{align*}
\hat{p}_E
&= \ln\bigl(p_{E,t+1}\bigr)  -  \ln\bigl(p_{E,t}\bigr)\\
&= \Bigl[\ln(p_G) + \ln\bigl(X_G + \tfrac{X_R}{\beta}\bigr) - \ln\bigl(\beta X_G + X_R\bigr)\Bigr]
    -  \Bigl[-\ln(p_G)
     - \bigl(\ln\bigl(X_G + \tfrac{X_R}{\beta}\bigr) - \ln(\beta X_G + X_R)\bigr)\Bigr]\\
&= 0.
\end{align*}
Hence, there is no change in the price of energy.

According to the dual definition of productivity growth, Eq. \ref{eq:dual_appendix}, since the price of factors are fixed,
\begin{equation}
    \hat{A}_E=0.
\end{equation}

NB: the cost share of the product G and R are not constant over time.

\subsubsection{Manufacturing M}

Since the price of energy and the quantity of final energy remain unchanged, we have
\begin{equation}
    \hat{A}_M = 0.
\end{equation}

\subsubsection{Aggregate TFP}

Thus aggregate productivity growth is
\begin{equation}
\hat{A} = 0.
\end{equation}


\subsection{Conclusion}

Under consistent input--output accounting (with fixed shares $a_{ji}$), partial substitution of gas-based electricity by renewables in manufacturing does not alter aggregate TFP growth, regardless of whether the energy flows are measured as final, useful, or primary.  Each method yields the same conclusion.


\section{With our framework}

In this subsection, we compute industry-level productivity growth as
\begin{equation}
    \hat{A}_i=\zeta \hat{X}_i+(\zeta-1)\hat{p}_i.
\end{equation}

\subsection{Case 1: Primary-Energy Approach}

For the gas, renewables industries, we have

\begin{equation}
    \begin{cases}
         \hat{X}_G^p&=\ln(1-\delta),\\
         \hat{p}_G&=0,\\
         \hat{A}_G &= \zeta \hat{X}_G^p=\zeta \ln(1-\delta),
    \end{cases}
 \quad\text{and}\quad
     \begin{cases}
         \hat{X}_R^p&=\ln(1+\gamma),\\
         \hat{p}_R&=0,\\
         \hat{A}_R &= \zeta \hat{X}_R^p=\zeta \ln(1+\gamma).
    \end{cases}
\end{equation}

And for the manufacturing sector, we have
\begin{equation*}
    \begin{cases}
         \hat{X}_M&=0,\\
         \hat{p}_M=&(\widehat{p_GX_G^p+p_RX_R^p})-\hat{X}_M=a_{GM} \ln(1-\delta) + a_{RM} \ln(1+\gamma))=0,\\
         \hat{A}_M &= 0.
    \end{cases}
\end{equation*}
Hence, applying the same method as in Eq. \ref{eq_final_Ahat_M},
\begin{equation}
    \hat{A} =
 \lambda_G \zeta\ln(1-\delta)
  + \lambda_R \zeta\ln(1+\gamma)
  =0.
\end{equation}


\subsection{Case 2: Final/Useful-Energy Approach}

For the gas and renewables industries, we have

\begin{equation}
    \begin{cases}
         \hat{X}_G^f&=\ln(1-\delta),\\
         \hat{p}_G&=0,\\
         \hat{A}_G &= \zeta \hat{X}_G^f=\zeta \ln(1-\delta).
    \end{cases}
    \quad\text{and}\quad
     \begin{cases}
         \hat{X}_R^f&=\ln(1+\gamma),\\
         \hat{p}_R&=0,\\
         \hat{A}_R &= \zeta \hat{X}_R^f=\zeta \ln(1+\gamma).
    \end{cases}
\end{equation}

Since the price of electricity remains the same regardless of whether it is generated from gas or renewables (Eq. \ref{eq:price_elec}), substituting one for the other in the manufacturing sector does not alter cost shares, input use, the final product’s price, or production levels. Consequently, the impact on productivity is null, 
\begin{equation*}
     \begin{cases}
         \hat{X}_M=&0,\\
         \hat{p}_M=&0,\\
         \hat{A}_M =& 0.
    \end{cases}
\end{equation*}
Hence, applying the same method as in Eq. \ref{eq_final_Ahat_M},
\begin{equation}
    \hat{A} =
 \lambda_G \zeta\ln(1-\delta)
  + \lambda_R \zeta\ln(1+\gamma)
  =0
\end{equation}


\subsection{Case 4: Primary energy with two sectors}

\paragraph{Energy sector}
In this case, $\delta X_G$ is being replaced by $\beta \delta X_G$, hence the production in MWh of the energy sector is no constant, however, sales $C_E$ are.


 Assuming $p_R=\frac{p_G}{\beta }$,
\begin{align*}
    \hat{C}_E&=\ln(p_G(1-\delta)X_G+p_R(X_R+X_G\delta \beta ))-\ln(p_GX_G+p_RX_R)\\
    &=\ln(p_GX_G+p_G\frac{X_R}{\beta })-\ln(p_GX_G+p_G\frac{X_R}{\beta })\\
&=0.
\end{align*}

Hence, since 
\begin{equation*}
    \hat{C}_E=\hat{X}_E^p+\hat{p}_E=0,
\end{equation*}
then,
\begin{equation}
   \hat{p}_E=-\hat{X}_E^p.  
\end{equation}

And productivity growth is
\begin{equation}
    \hat{A}_E=\zeta \hat{X}_E^p+(\zeta-1)\hat{p}_E=\hat{X}_E^p
\end{equation}

\paragraph{Manufacturing}

Productivity growth of the manufacturing sector,
\begin{equation*}
    \hat{A}_M=\zeta\hat{X}_M+(\zeta-1)\hat{p}_M
\end{equation*}

By assumption, output growth =is $\hat{X}_M=0$, and due to the null profit assumption, the costs of the M sector are equal to sales, $p_EX_E^p+L_fw_f=p_MX_M$. Thus, we compute price growth as
\begin{align*}
    \hat{p}_M&=\widehat{\left(\frac{p_EX_E^p+L_fw_f}{X_M}\right)}\\
    &=\widehat{p_EX_E^p+L_fw_f}-\underbrace{\hat{X}_M}_0\\
     &=\widehat{p_EX_E^p+L_fw_f}\\
    &=\frac{p_EX_E}{p_EX_E^p+L_fw_f}\underbrace{\widehat{p_EX_E^p}}_{0}+\frac{L_fw_f}{p_EX_E^p+L_fw_f}\underbrace{\widehat{L_fw_f}}_{0}\\
    &=0.
\end{align*}
Thus, 
\begin{equation}
    \hat{A}_M=0.
\end{equation}

\paragraph{Aggregate TFP}
Aggregate productivity growth is

\begin{align*}
        \hat{A}= \lambda_E \hat{X}_E^p=\frac{p_EX_E^p}{Y}\hat{X}_E^p
\end{align*}



\subsection{Case 5: Final-Energy (resp: Useful) Approach for two industries}
\label{subappendix_final_energy}


\subsubsection{Energy industry}

In final electricity or useful energy approach, the energy demand from M is constant through the electrification switch.

The energy sector production growth is
\begin{equation*}
   \hat{X}_E^f= \ln(\beta X_G(1-\delta)+(X_R+\beta X_G\delta))-\ln(\beta X_G+X_R)=0
\end{equation*}
Likewise, 
\begin{equation*}
\hat{p}_E=  \ln(p_G) +\ln\left[X_G+\frac{X_R}{\beta }\right] - \ln(\beta X_G+X_R)-
\ln(p_G)-\left(\ln(X_G +\frac{X_R}{\beta } ) - \ln(\beta X_G+X_R)\right)=  0.
\end{equation*}
Hence, 
\begin{equation}
    \hat{A}_E=0
\end{equation}

In the manufacturing sector, the price of energy and quantity of useful do not change, hence neither the production nor the price of manufacturing output chances. Therefore
\begin{equation}
    \hat{A}_M=0
\end{equation}

Aggregate productivity growth is
\footnote{
Even if we considered that here we had introduced two new inputs: G and R. 
For each of these sectors, we have
\begin{align*}
    \hat{A}_G&=\zeta \ln(1-\delta)\\
    \hat{A}_R&=\zeta \ln(\gamma)
\end{align*}
Thus
\begin{equation*}
\hat{A} = \lambda_G\hat{A}_G+\lambda_R\hat{A}_R=0
\end{equation*}
}

\begin{equation}
\hat{A} =0
\end{equation}




\subsection{Draft}

\subsubsection{Other approach - primary energy, two sectors, national accounting}
NB: we have actually introduced two new inputs: G and R.
So dual of Energy doesn't change, but primal
\begin{align*}
    \hat{A}_E&=\hat{X}_E-a_{GE}\hat{X}_G-a_{RE}\hat{X}_R \\
    &=\frac{X_G}{X_G+X_R} \hat{X}_G+\frac{X_R}{X_G+X_R} \hat{X}_E-\frac{p_GX_G}{p_EX_E}\hat{X}_G-\frac{p_RX_R}{p_EX_E}\hat{X}_R\\
    &= \left(\frac{X_G}{X_E} -\frac{p_GX_G}{p_EX_E}\right)\hat{X}_G+\left(\frac{X_R}{X_E} -\frac{p_RX_R}{p_EX_E}\right)\hat{X}_E\\
      &= \left(\frac{p_EX_G-p_GX_G}{p_EX_E} \right)\hat{X}_G+\left(\frac{p_EX_R-p_RX_R}{p_EX_E}\right)\hat{X}_E\\
\end{align*}

Agg
\begin{align*}
    \hat{A}&=\lambda_E \hat{A}_E+\lambda_M \hat{A}_M+\lambda_G\hat{A}_G+\lambda_R\hat{A}_R\\
    &=\lambda_E \hat{A}_E+\lambda_M \hat{A}_M\\
    &=\frac{p_EX_E}{Y}\left[\left(\frac{p_EX_G-p_GX_G}{p_EX_E} \right)\hat{X}_G+\left(\frac{p_EX_R-p_RX_R}{p_EX_E}\right)\hat{X}_E\right]+\frac{p_XX_M}{Y} (-a_{EM}\hat{X}_E)\\
     &=\left[\left(\frac{p_EX_G-p_GX_G}{Y} \right)\hat{X}_G+\left(\frac{p_EX_R-p_RX_R}{Y}\right)\hat{X}_E\right]+\frac{X_Ep_E}{Y}\hat{X}_E\\
\end{align*}


\subsubsection{National Accounting : Case 5, final energy, two sectors}


Let us check whether the share of factor are constant;

Before the switch,
\begin{align*}
    a_{GE,t}=\frac{p_G X_G}{p_G X_G+p_R X_R}=\frac{X_G}{X_G+\frac{X_R}{\beta }}.
\end{align*}

After the switch,
\begin{align*}
    a_{GE,t+1}&=\frac{p_G (1-\delta)X_G}{p_G (1-\delta)X_G+p_R (X_R+X_G\delta \beta )}\\
    &=\frac{p_G (1-\delta)X_G}{p_G (1-\delta)X_G+\frac{p_G}{\beta } (X_R+X_G\delta \beta )}\\
    &=\frac{(1-\delta)X_G}{(1-\delta)X_G+\frac{1}{\beta } (X_R+X_G\delta \beta )}\\
    &=\frac{(1-\delta)X_G}{(1-\delta)X_G+\frac{X_R}{\beta } +X_G\delta}\\
     &=\frac{(1-\delta)X_G}{X_G+\frac{X_R}{\beta }}\\
     &=(1-\delta)a_{GE}.
\end{align*}

\begin{align*}
    a_{RE}=\frac{p_R X_R}{p_G X_G+p_R X_R}=\frac{\frac{X_R}{\beta }}{X_G+\frac{X_R}{\beta }}
\end{align*}

\begin{align*}
    a_{RE,t+1}&=\frac{p_R (X_R+X_G\delta \beta )}{p_G (1-\delta)X_G+p_R (X_R+X_G\delta \beta )}\\
    &=\frac{\frac{p_G}{\beta } (X_R+X_G\delta \beta )}{p_G (1-\delta)X_G+\frac{p_G}{\beta } (X_R+X_G\delta \beta )}\\
    &=\frac{\frac{1}{\beta }X_R + X_G\delta}{(1-\delta)X_G+\frac{1}{\beta } (X_R+X_G\delta \beta )}\\
    &=\frac{\frac{1}{\beta }X_R + X_G\delta}{(1-\delta)X_G+\frac{X_R}{\beta } +X_G\delta}\\
     &=\frac{\frac{1}{\beta }X_R + X_G\delta}{X_G+\frac{X_R}{\beta }}\\
    &=\frac{\delta(X_G+\frac{X_R}{\beta })+\frac{X_R}{\beta }-\delta\frac{X_R}{\beta }}{X_G+\frac{X_R}{\beta }}\\
    &=\frac{\delta(X_G+\frac{X_R}{\beta })+\frac{X_R}{\beta }(1-\delta)}{X_G+\frac{X_R}{\beta }}\\
     &=\delta +(1-\delta)a_{RE}
\end{align*}

The factor shares are not constant. We could do Divisia/Tornqvist (discrete time and take the average) or just use the first period ones.

\begin{equation*}
  \begin{cases}
       \hat{X}_G&=\ln(1-\delta)\\
        \hat{X}_R&=\ln(X_R+\beta \delta X_G)-\ln(X_R)=\ln(\gamma)
  \end{cases}
\end{equation*}

\begin{equation*}
    \hat{A}_E=-\hat{X}_G \frac{X_G}{X_G+\frac{X_R}{\beta }}-\hat{X}_R \frac{\frac{X_R}{\beta }}{X_G+\frac{X_R}{\beta }}
\end{equation*}

\subsubsection{Framework, Case Primary, two sectors}



\paragraph{Energy}

\begin{equation*}
    \hat{A}_E=\zeta\hat{X}_E+(\zeta-1)\hat{p}_E
\end{equation*}

X
\begin{align*}
    \hat{X}_E=
&= \widehat{X_G+X_R} \\
&= \frac{X_G}{X_G+X_R} \hat{X}_G+\frac{X_R}{X_G+X_R} \hat{X}_R\\
&=\frac{X_G}{X_G+X_R} (\ln(X_G(1-\delta)-\ln(X_G))+\frac{X_R}{X_G+X_R} (\ln(X_R+X_G\beta \delta)-\ln(X_R))\\
&=\frac{X_G}{X_G+X_R} \ln(1-\delta)+\frac{X_R}{X_G+X_R} \ln\left(1+\frac{X_G}{X_R}\beta \delta\right)\\
\end{align*}



Average price: 
\begin{align*}
    \ln(p_{E,t+1})&=
    \ln\frac{X_{G,t+1}p_G +X_{R,t+1}p_R}{X_{G,t+1} +X_{R,t+1}}\\&=
    \ln \left(X_{G,t+1}p_G +X_{R,t+1}p_R \right)-\ln\left(X_{G,t+1} +X_{R,t+1}\right)\\
    &=
    \left(\ln\left[(1-\delta)X_G p_G+(X_R+\delta X_G \beta ) p_R\right] - \ln((1-\delta)X_G+(X_R+X_G \delta \beta ) )\right)\\
\end{align*}

\begin{align*}
    \ln(p_{E,t})&=
    \ln\frac{X_{G,t}p_G +X_{R,t}p_R}{X_{G,t} +X_{R,t}}\\&=
    \ln \left(X_{G,t}p_G +X_{R,t}p_R \right)-\ln\left(X_{G,t} +X_{R,t}\right)\\
    &=
    \ln(X_G p_G+X_R p_R) - \ln(X_G+X_R)\\
\end{align*}


\begin{align*}
    \hat{p}_E&=\ln(p_{E,t+1})-\ln(p_{E,t})\\
    &= \ln\left[(1-\delta)X_G p_G+(X_R+\delta X_G \beta ) p_R\right] - \ln((1-\delta)X_G+(X_R+X_G \delta \beta ) )
    -
    \ln(X_G p_G+X_R p_R) + \ln(X_G+X_R)\\
     &\quad \text{using}\quad p_R=\frac{p_G}{\beta },\\
     &=\ln\left[(1-\delta)X_G p_G+\frac{p_G}{\beta }X_R+\delta X_Gp_G \right] - \ln((1-\delta)X_G+(X_R+X_G \delta \beta ) )
    -
    \ln(X_G p_G+X_R \frac{p_G}{\beta }) + \ln(X_G+X_R)\\
     &=\ln(p_G)+\ln\left[X_G+\frac{X_R}{\beta } \right] - \ln((1-\delta)X_G+(X_R+X_G \delta \beta ) )-\ln(p_G)-    \ln(X_G+ \frac{X_R}{\beta }) + \ln(X_G+X_R)\\
     &=- \ln((1-\delta)X_G+(X_R+X_G \delta \beta ) )+\ln(X_G+X_R)\\
       &=- \ln(X_G+X_R-\delta X_G(1+\beta )) )+\ln(X_G+X_R)\\
           &=- \ln\left(1-\delta(1+\beta )\frac{ X_G}{X_G+X_R}\right)\\
\end{align*}

Pour rappel, 
\begin{align*}
    \hat{X}_E&=\frac{X_G}{X_G+X_R} \ln(1-\delta)+\frac{X_R}{X_G+X_R} \ln\left(1+\frac{X_G}{X_R}\beta \delta\right)\\
\end{align*}


\subsubsection{Case 4: Primary with two sectors}
The switched quantity of gas-based electricity, is exactly equal to the additional amount of renewable electricity,
\begin{equation*}
\label{eq_constant_final_demand_elec}
    \Delta X_R = \delta X_G.
\end{equation*}


\paragraph{Energy industry E} 

\begin{equation*}
    \delta q_p^G \beta _{fG}= \Delta X_R
\end{equation*}

\begin{equation*}
\hat{p}_E=\hat{C}_E-\hat{X}_E   
\end{equation*}


\begin{align*}
    \hat{C}_E&=\ln(p_G(1-\delta)X_G+p_R(X_R+X_G\delta \beta ))-\ln(p_GX_G+p_RX_R)\\
    &=\ln(p_GX_G+p_G\frac{X_R}{\beta })-\ln(p_GX_G+p_G\frac{X_R}{\beta })\\
&=0
\end{align*}

\begin{align*}
    \hat{C}_E&=\ln(p_G(1-\delta)X_G+p_R(X_R+X_G\delta \beta ))-\ln(p_GX_G+p_RX_R)\\
    &=\ln(p_GX_G+p_G\frac{X_R}{\beta })-\ln(p_GX_G+p_G\frac{X_R}{\beta })\\
&=0
\end{align*}

Hence
\begin{equation}
   \hat{p}_E=-\hat{X}_E    
\end{equation}

\begin{equation*}
\begin{aligned}
    \hat{X}_E
&= \widehat{X_G+X_R} \\
&= \frac{X_G}{X_G+X_R} \hat{X}_G+\frac{X_R}{X_G+X_R} \hat{X}_R\\
&= \frac{X_G}{X_G+X_R} \ln(1-\delta)+\frac{X_R}{X_G+X_R} \ln(\gamma)\\
\end{aligned}
\end{equation*}

\begin{align*}
    \hat{A}_E& = \zeta \hat{X}_E+ (\zeta-1)\hat{p}_E\\
    &=\zeta \hat{X}_E- (\zeta-1)\hat{X}_E\\
    & = \hat{X}_E  =-\hat{p}_E\\
\end{align*}


\paragraph{Manufacturing M}

\begin{equation*}
    \hat{X}_M=0
\end{equation*}

price
\begin{align*}
    \hat{p}_M&=\widehat{\left(\frac{C_E+C_F}{X_M}\right)}\\
    &=\widehat{C_E+C_F}-\underbrace{\hat{X}_M}_0\\
     &=\widehat{C_E+C_F}\\
    &=\frac{C_E}{C_E+C_F}\underbrace{\hat{C}_E}_0+\frac{C_F}{C_E+C_F}\underbrace{\hat{C}_F}_0\\
    &=0
\end{align*}

Hence, 
\begin{equation}
    \hat{A}_M=0
\end{equation}


\paragraph{Aggregate} 


But here we have introduced two new inputs: G and R. 
For each of these sectors: 
\begin{align*}
    \hat{A}_G&=\zeta \ln(1-\delta)\\
    \hat{A}_R&=\zeta \ln(\gamma)
\end{align*}


\begin{align*}
        \hat{A}&=\lambda_E \hat{X}_E+\lambda_G \zeta\hat{X}_G+\lambda_R \zeta\hat{X}_R \\
        &= \frac{p_EX_E}{Y}
        \left[
        \hat{X}_E+\zeta \left(\frac{p_GX_G}{p_EX_E} \hat{X}_G+\frac{p_RX_R}{p_EX_E} \hat{X}_R\right)
        \right]\\
        &= \frac{p_EX_E}{Y}
        \left[
        \hat{X}_E+\zeta \left(\frac{p_GX_G}{p_GX_G+p_RX_R} (\hat{X}_G+\hat{p}_G)+\frac{p_RX_R}{p_GX_G+p_RX_R} (\hat{X}_R+\hat{p}_R\right)
        \right]\\
        &= \frac{p_EX_E}{Y}
        \left[
        \hat{X}_E+\zeta \widehat{(C_R+C_G)}
        \right]\\
        &= \frac{p_EX_E}{Y}  \hat{X}_E\\
        &=\lambda_E \hat{X}_E
\end{align*}
% \begin{align*}
%     \hat{A}&=\lambda_E \hat{A}_E+\lambda_M \hat{A}_M\\
%     &=\frac{X_Ep_E}{Y} (\zeta \hat{X}_E+ (\zeta-1)\hat{p}_E)+ \frac{C_E+C_F}{Y} (\zeta-1)\frac{C_E}{C_E+C_F}\left(\hat{X}_E+\hat{p}_E\right)\\
%     &=\frac{C_E}{Y} (\zeta \hat{X}_E+ (\zeta-1)\hat{p}_E)+ \frac{C_E}{Y} (\zeta-1)\left(\hat{X}_E+\hat{p}_E\right)\\
%      &=\hat{X}_E \left[\zeta   \frac{C_E}{Y} +\frac{C_E}{Y} (\zeta-1) \right]+
%      \hat{p}_E \left[(\zeta-1)   \frac{C_E}{Y} +\frac{C_E}{Y} (\zeta-1) \right]
% \end{align*}

Note that even without the assumption $p_R=\frac{p_G}{\beta }$\footnote{It still checks out with the previous case if $\hat{p}_E=-\hat{X}_E$}, we would have the same result for aggregate productivity growth
\begin{align*}
        \hat{A}&=\lambda_E (\zeta\hat{X}_E+(\zeta-1)\hat{p}_E)+\lambda_M (\zeta-1)\frac{C_E}{C_E+C_F}\hat{C}_E\\
        &= \frac{C_E}{Y}
        \left[\zeta\hat{X}_E+(\zeta-1)\hat{p}_E)\right]+
        \frac{C_E+C_F}{Y}(\zeta-1)\frac{C_E}{C_E+C_F}(\hat{X}_E+\hat{p}_E)\\
        &= \frac{C_E}{Y}
        \left[(2\zeta-1)\hat{X}_E+2(\zeta-1)\hat{p}_E)\right].\\
\end{align*}

